\chapter*{Conclusion}
\label{conclusion}
\addcontentsline{toc}{chapter}{Conclusion}
\markboth{Conclusion}{Conclusion}

Au fil de ce livre, nous avons parcouru toutes les étapes d'une enquête avec Epi Info 7 : concevoir un questionnaire robuste, le sécuriser avec le code de contrôle, saisir les données, puis les analyser avec le \module{tableau de bord visuel} et le module \module{classique}. Avec \module{statcalc}, nous avons aussi vu comment calculer la taille d'un échantillon avant même de commencer la collecte.

Si vous ne deviez retenir que quelques principes, ce seraient ceux-ci :
\begin{itemize}[label=\ding{97}, font=\color{EPIB}]
	\item une bonne analyse commence par un bon questionnaire : des noms de variables clairs, les bons types de champs et des contrôles de saisie ;
	\item testez toujours le questionnaire avant la collecte, puis ne modifiez plus sa structure ;
	\item vérifiez l'effectif de chaque analyse, surtout lorsqu'un filtre ou une sélection est actif ;
	\item enregistrez vos programmes d'analyse : ils rendent votre travail reproductible ;
	\item une association statistiquement significative ne prouve pas un lien de causalité.
\end{itemize}

Epi Info n'est qu'un outil. Ce qui donne de la valeur à vos résultats, c'est la rigueur de votre démarche, de la conception du protocole à l'interprétation des chiffres. Une fois ces bases acquises, vous pourrez aborder sereinement des analyses plus avancées, ou d'autres logiciels comme \textsc{r} ou \textsc{stata}.

Je vous souhaite de belles enquêtes et de bonnes analyses.

\bigskip
\begin{flushright}
	\textit{OUEDRAOGO Ahmed Dieudonné}\\
	\textit{2026}
\end{flushright}
